\section{Actual original finite tensor forms retain their whole joint-chart coefficients and generic values}
\begin{lemma}\label{lem:original-subopen-frame-coefficient-of-scalar}
\lean{EnvelopeProperty.OriginalProjectiveTensorPivotCoefficients.frame_coefficient_of_scalar}
\leanok
\uses{thm:actual-pullback-subopen-scalar,sc0-native-line-ratio:whole-frame,lem:fc-original-local-section-scalar}
A genuine whole original scalar equation in an actual pulled frame on an original affine open computes the actual native local-trivialization coefficient of that same original global section. The frame on the restricted module is the canonical restriction comparison of the same pulled frame, and the section is the value at one of the same original section morphism.
\end{lemma}
\begin{proof}
The genuine section-to-morphism and value-at-one maps are inverses. Restriction and pullback frames compare by the original canonical module isomorphism. Apply the whole local scalar identity of the actual section in its restricted frame. The assumed genuine scalar equation is an equality of the same original module morphisms. Evaluate at one, then use injectivity of the actual structure-section top-open isomorphism to identify the original coefficients.
\end{proof}
\begin{definition}\label{def:original-finite-projective-tensor-coefficients}
\lean{EnvelopeProperty.OriginalProjectiveTensorPivotCoefficients.chartFrame,EnvelopeProperty.OriginalProjectiveTensorPivotCoefficients.coefficient,EnvelopeProperty.OriginalProjectiveTensorPivotCoefficients.formSection}
\leanok
\uses{thm:finite-original-projective-power-joint-frames}
For the actual finite original tensor of pulled projective powers, use the actual joint frame on any original affine open contained in all coordinate preimages. Its coefficient map is the genuine local-trivialization coefficient of every actual original whole section. The original tuple of homogeneous forms gives its same whole tensor section.
\end{definition}
\begin{theorem}\label{thm:original-finite-projective-tensor-form-coefficients}
\lean{EnvelopeProperty.OriginalProjectiveTensorPivotCoefficients.coefficient_formSection}
\leanok
\uses{def:original-finite-projective-tensor-coefficients,lem:original-subopen-frame-coefficient-of-scalar,thm:finite-original-projective-power-joint-frames}
The actual original whole tensor form section has, in this SAME joint frame, the product of every original dehomogenized coefficient pulled through its genuine restricted factor map. This is an identity of actual original regular functions on the joint affine open. No form, scalar or frame matching certificate is supplied.
\end{theorem}
\begin{proof}
The actual finite pulled tensor-frame scalar theorem gives the product in the true top-open section ring. The original affine pullback-section construction identifies each factor with its original ambient regular function on this same joint open. The top-open structure isomorphism preserves the finite product. The preceding actual scalar-to-coefficient lemma proves the asserted original coefficient equality.
\end{proof}
\begin{theorem}\label{thm:original-finite-projective-tensor-form-generic-coefficients}
\lean{EnvelopeProperty.OriginalProjectiveTensorPivotCoefficients.germ_coefficient_formSection}
\leanok
\uses{thm:original-finite-projective-tensor-form-coefficients,lem:fc-original-homogeneous-form-frame-coefficient}
On the SAME original integral scheme, the generic germ of every such joint-chart coefficient is the product of the original homogeneous forms evaluated at the SAME actual original projective ratios for those coordinate charts. It is an equality in the actual original function field. Only nonemptiness of the original joint open is needed, and nonemptiness of each original factor chart is derived from that inclusion.
\end{theorem}
\begin{proof}
The original structure-ring germ preserves the finite product. Each genuine affine pullback equals the original ambient section comorphism followed by its original restriction to the joint open. Restriction to a nonempty original open preserves the same generic germ. The original homogeneous coefficient theorem computes that original germ by actual polynomial evaluation at the original projective coordinate ratios. Apply it to every original factor, retaining its actual original coefficient action.
\end{proof}
\begin{theorem}\label{thm:original-finite-projective-tensor-pivot-nonzero}
\lean{EnvelopeProperty.OriginalProjectiveTensorPivotCoefficients.pivotSection,EnvelopeProperty.OriginalProjectiveTensorPivotCoefficients.germ_coefficient_pivotSection,EnvelopeProperty.OriginalProjectiveTensorPivotCoefficients.pivotSection_ne_zero}
\leanok
\uses{thm:original-finite-projective-tensor-form-generic-coefficients}
On the same original integral scheme, choose one original nonempty coordinate chart of each projective factor. The tuple of original coordinate powers constructs a whole original tensor section. Its generic coefficient in every original joint frame is the product of the corresponding true original coordinate ratios to the original powers. This constructed whole pivot tensor section is nonzero, including when some or all original degrees vanish. No nonzero section witness is supplied.
\end{theorem}
\begin{proof}
Apply the preceding whole form coefficient formula to each original coordinate power. Actual polynomial evaluation gives the stated powers of the original ratios. Every chosen nonempty factor chart contains the original generic point, so their finite intersection also contains it. Choose an original affine neighborhood of that same generic point inside the intersection. In its joint pivot frame, every original pivot ratio is one, so the generic coefficient of the constructed whole tensor section is one. The zero section has zero original coefficient and zero original generic germ, contradicting that equality. This proves actual nonvanishing from the original geometry.
\end{proof}
Whole-section gluing and norm comparison remain the next derived steps. Arbitrary signed multidegrees will use actual dual ratios rather than an assumed nonzero signed global section. Whole projective numerical finite-dimensionality and the paper topology remain open.
